\scanpage{279}


Ainsi on peut écrire

\[
u'=u+\unsure{\zeta|\xi}=u-\zeta\qquad(\text{car }2\zeta=0)
\]

où $\zeta$ est un des trois éléments $\ne0$ de $2M\simeq M\otimes_\Lambda k_{i_0}$, et plus précisément, un des deux éléments \deleted{\unsure{$\ne$}}

\[
\zeta\in[(2M=M\otimes_\Lambda k_{i_0})\setminus\{0,2u\}]
\]

(NB $\zeta\notin\{0,2u\}$ signifie justement que $u+\zeta\notin\{\pm u\}$)

Si

\[
2M\setminus\{0,\underset{\unsure{\eta|?}}{\underbrace{2u}}\}=\{\zeta',\zeta''\},
\]

notons que \unsure{l'on a|alors}

\[
\deleted{\unsure{2u}}\underset{2u}{\underbrace{\eta}}+\zeta'+\zeta''=0
\]

donc

\[
u+\zeta'=\underset{-u-\zeta''}{\underbrace{-(u+\zeta'')}}\quad\unsure{\text{i.e.}}
\]

i.e. l'image de $u'$ dans $M^{\unsure{*}}/\{\pm\operatorname{id}_M\}=S_M$ est uniquement déterminée par $u$. \deleted{Ainsi} On \deleted{peut} a ainsi une bijection canonique

\[
\underline a:S_M\longrightarrow S_M
\]

qui associe à tout $\dot u\in S_M$ l'unique $\dot u'\in S_M$ tel que l'on ait

\[
\operatorname{Et}^{\deleted{\unsure{*}}}_{i}(\dot u)=\operatorname{Et}_i(\dot u'),\quad \dot u'\ne\dot u
\]

pour un $i\in\unsure{\Pi_M|\pi_M}=(\det M)^*/\pm1_\Lambda$ -- ce qui est alors vrai pour tous. On voit de suite que $\underline a$ est involutif, \inserted{et} ss pts fixes par construction.

On peut voir \deleted{\unsure{cet|les}}\inserted{décrire} antipodisme par la suite exacte

\[
0\longrightarrow\deleted{\unsure{\operatorname{End}(M_0)}}\underset{\simeq\operatorname{End}_{\mathbb F_2}(M_0)}{\underbrace{\unsure{N|\mathcal N}}}\longrightarrow\operatorname{Aut}_\Lambda(M)\longrightarrow\operatorname{Aut}_{\Lambda/2\Lambda}(M/2M)\longrightarrow1
\]

où $M_0=M\otimes_\Lambda(k_{i_0}\simeq\mathbb F_2)$, et on a

\[
\begin{aligned}
&\operatorname{End}(M_0)\simeq M_0\oplus\mathbb F_2^{\unsure{\underline\omega_J|\omega_J|\simeq J}}\qquad(J=M_0^*),\\
\text{Donc}\quad&\deleted{\unsure{\operatorname{End}(M_0)}}N/\{\pm\operatorname{id}_M\}\simeq M_0\oplus\mathbb F_2^{\underline\omega_J}/\text{diag}\simeq M_0+\mathbb F_2
\end{aligned}
\]

et l'élément non nul du facteur $\mathbb F_2$ donne $\underline a$.
% archive sheet number: 278
